\section{Conclusions}
\label{sec:conclusion}

\engiproof{} converts selected published engineering methods into source-bounded evidence. It keeps separate what was published, what was reproduced, what was checked independently, what disagrees, what a human decided and what remains unqualified. The six heterogeneous studies support four statements:
\begin{itemize}[leftmargin=1.5em]
  \item \textbf{Reproduction.} Bounded targets are reproduced to source precision, and the outcome is recorded target by target.
  \item \textbf{Independent checks.} Methodologically independent checks decide the direction of numerical conflicts without tuning.
  \item \textbf{Discrepancies.} Source inconsistencies and missing inputs are preserved as evidence, and they gate promotion until a human decides.
  \item \textbf{Verification.} Verification recomputes in isolation and classifies differences across operating systems without mutating the evidence it checks.
\end{itemize}

The strongest conclusion is methodological. Engineering evidence automation should make disagreement and incompleteness explicit and traceable. A paper-derived method should not be promoted because its code runs, or because a printed number can be made to reappear.

Before submission, four owner actions remain:
\begin{itemize}[leftmargin=1.5em]
  \item record the pending reviewer decisions on discrepancies;
  \item approve or amend the proposed taxonomy mapping;
  \item approve or amend the cross-environment comparator rules;
  \item generate the remaining ingestion summaries from the local intakes.
\end{itemize}
The claim--evidence matrix will then be regenerated from the tagged repository state.
